% SI-ready method fragment. Keep this fragment separate until the coauthor
% manuscript review explicitly decides where it enters sn-article-SI.tex.
\subsection{Failure-attribution candidate classification and selection}
\label{sec:si_candidate_selection}

We constructed a provenance-preserving inventory from the prepared-only
snapshots of the all-replicate failure-attribution analysis. Each candidate was
a saved 2050-atom fully periodic configuration. Frame-level diagnostics were
joined using the exact tuple $(\texttt{source},\texttt{training\_seed},
\texttt{window},\texttt{step})$, and run-level metadata using
$(\texttt{source},\texttt{training\_seed})$. No frame was discarded because it
was unstable. After deterministic sorting by source, training seed, step, and
window order, we assigned canonical identifiers
\texttt{dftcand\_0001}--\texttt{dftcand\_0639}.

The source windows were relabeled as follows: \texttt{stable\_early} became
\texttt{stable\_early\_control}, \texttt{stable\_pre} became
\texttt{stable\_pre\_onset}, and \texttt{precursor} became
\texttt{precursor\_near\_onset}. A threshold frame with an observed trigger
was labeled \texttt{threshold\_or\_terminal}; a threshold window with
\texttt{trigger = not\_observed} was labeled
\texttt{terminal\_stable\_target}. The resulting inventory contained 207,
216, 208, 7, and 1 frames in these five groups, respectively.

The original 20-frame coverage wave retained all seven observed
threshold/terminal rows and the one 1000~ps stable endpoint, selected the latest
sampled precursor frame from each of the nine source/training-seed runs, and
added one approximately 10~ps stable-early control for each training seed.
This wave therefore contained 7 threshold, 1 stable endpoint, 9 precursor,
and 3 early-stable controls. The exact row-level ledger is archived with the
reproduction package.

For the subsequent MLIP-informed DFT prioritization, every candidate was
evaluated as a static single point with three SHA256-identified checkpoints.
We first required successful results from all three models and then imposed a
source-geometry gate of $d_{\min}\geq0.8$~\AA. Within each source-stage group,
we ranked candidates by decreasing ensemble force-spread RMS, decreasing
maximum pairwise force difference, and increasing canonical ID. We applied
quotas of 4 threshold/terminal, 8 precursor, 4 stable-pre, 2 early-stable,
and 1 stable-terminal candidate; if the geometry gate left a quota unfilled,
the remaining slots were filled from unused gate-passing rows in the same
priority order. The realized 20-frame subset contained 1 threshold/terminal,
12 precursor, 4 stable-pre, 2 early-stable, and 1 stable-terminal candidate.

The three checkpoints were the $l_{\max}=2$ ablation models trained with seeds
123 and 234 and the baseline model trained with seed 345. The source MD
velocity seed (20261005) was recorded separately from the training seed. The
static force-spread score is used only to prioritize independent reference
calculations; it is not a calibrated uncertainty and does not establish that
any model force is incorrect or causal. The reduced-cell QE inputs are
explicitly reconstructions of the full-cell periodic snapshots, so they are
not promoted to full-cell DFT parity until the corresponding convergence and
identity gates pass.
